\documentclass[aspectratio=169,10pt]{beamer}

\usepackage{iftex}
\ifPDFTeX
  \usepackage[T1]{fontenc}
  \usepackage[utf8]{inputenc}
  \usepackage{lmodern}
\fi
\usepackage[spanish,es-noshorthands]{babel}
\usepackage{booktabs}
\usepackage{array}
\usepackage{graphicx}
\usepackage{tikz}
\usetikzlibrary{arrows.meta,positioning,shapes.geometric,calc,babel}

% Paleta común para todos los diagramas
\colorlet{cSensor}{blue!60!black}
\colorlet{cProc}{orange!80!black}
\colorlet{cAct}{green!45!black}
\colorlet{cApp}{violet!70!black}
\colorlet{cPower}{gray!70!black}

\tikzset{
  >={Stealth[length=2mm]},
  flecha/.style={->, thick, draw=black!75},
  caja/.style={rectangle, rounded corners=2pt, draw=#1, fill=#1!10, thick,
    minimum width=2.3cm, minimum height=0.85cm, align=center, inner sep=2pt},
  etiqueta/.style={font=\scriptsize, text=black!70, fill=white, inner sep=1pt},
}


% ---------------------------------------------------------------
% Tema académico sobrio
\usetheme{Boadilla}
\definecolor{azul}{HTML}{1F3A5F}
\definecolor{gris}{HTML}{5F6B7A}
\usecolortheme[named=azul]{structure}
\setbeamercolor{block title}{bg=azul!12, fg=azul}
\setbeamercolor{block body}{bg=azul!4}
\setbeamercolor{alerted text}{fg=cProc}
\setbeamertemplate{blocks}[rounded][shadow=false]
\setbeamertemplate{navigation symbols}{}
\setbeamertemplate{itemize items}[circle]
\setbeamertemplate{bibliography item}[text]
\setbeamerfont{frametitle}{series=\bfseries}

% Estilos TikZ propios de las diapositivas (fuera de los frames por el #)
\tikzset{
  f/.style={caja=#1, minimum width=1.85cm, minimum height=0.95cm},
  d/.style={diamond, aspect=1.5, draw=cProc, fill=cProc!10, thick, align=center,
    inner sep=0pt, minimum width=1.9cm, font=\scriptsize},
  capa/.style={rectangle, minimum width=2.4cm, minimum height=0.9cm,
    fill=#1, text=white, font=\small\bfseries},
  nuestro/.style={rectangle, draw=#1, fill=#1!8, thick, minimum height=0.9cm,
    align=left, font=\small, text width=3.2cm},
  n/.style={caja=#1, minimum width=1.8cm, minimum height=1cm},
}

% Foto de referencia con leyenda: \foto{archivo}{alto}{título}{detalle}
\newcommand{\foto}[4]{%
  \begin{minipage}[t]{\linewidth}\centering
    \fcolorbox{gris!40}{white}{\begin{minipage}[c][#2][c]{0.92\linewidth}\centering
      \includegraphics[height=0.94\dimexpr#2\relax, width=0.94\linewidth, keepaspectratio]{figuras/fotos/#1}
    \end{minipage}}\\[1mm]
    {\footnotesize\bfseries #3}\\[-0.5mm]
    {\scriptsize\color{gris} #4}
  \end{minipage}}
\setlength{\fboxsep}{0pt}
\setlength{\fboxrule}{0.3pt}

% Crédito de imágenes, sobre el pie de página
\newcommand{\credito}[1]{%
  \begin{tikzpicture}[remember picture, overlay]
    \node[anchor=south west, font=\tiny, text=gris, inner sep=0pt, text width=15cm]
      at ([xshift=4mm, yshift=6.5mm]current page.south west) {#1};
  \end{tikzpicture}}

\title[Seguridad residencial IoT]{Sistema inteligente de seguridad residencial IoT\\
  con monitoreo inalámbrico y alertas en tiempo real}
\subtitle{Perfil de proyecto (P1)}
\author[Salcedo, Quispe, Iribar, Salazar]{Ariann Abigail Salcedo Quiroz \and
  Diego Randolp Quispe Apaza \\ Federico Iribar Casanova \and
  Sofia del Pilar Salazar Mendoza}
\institute[UTEC]{Internet de las Cosas\\Universidad de Ingeniería y Tecnología}
\date[Set. 2026]{Setiembre de 2026}

\begin{document}

\begin{frame}
  \titlepage
\end{frame}

% ===============================================================
\section{Introducción}

\begin{frame}{Introducción}
  \begin{itemize}
    \item Una vivienda suele protegerse con piezas que no se hablan: cerradura, alarma y, con suerte, una cámara.
    \item Cada una trabaja sola. La casa tarda en notar un riesgo y casi nunca reacciona por sí misma.
    \item Hoy un ESP32 de pocos dólares puede leer sensores, decidir y avisar en tiempo real \cite{alfuqaha2015survey, oktafiani2025esp32}.
  \end{itemize}
  \vspace{4mm}
  \begin{block}{Propuesta}
    \centering
    \begin{tikzpicture}[font=\small, node distance=6mm]
      \node[caja=cSensor, minimum width=2.9cm] (a) {Detectar\\{\scriptsize PIR, magnético, vibración}};
      \node[caja=cProc, minimum width=2.9cm, right=of a] (b) {Decidir\\{\scriptsize ESP32, ventana de 10\,s}};
      \node[caja=cAct, minimum width=2.9cm, right=of b] (c) {Responder\\{\scriptsize sirena y cerradura}};
      \node[caja=cApp, minimum width=2.9cm, right=of c] (d) {Avisar y medir\\{\scriptsize app web en AWS}};
      \foreach \u/\v in {a/b, b/c, c/d} \draw[flecha] (\u) -- (\v);
    \end{tikzpicture}
  \end{block}
\end{frame}

% ===============================================================
\section{Problemática}

\begin{frame}{Problemática}
  Una intrusión puede empezar abriendo una puerta, con movimiento en una zona vacía o con golpes contra una estructura. Los sistemas recientes combinan PIR y contactos magnéticos para alertar apenas empieza el acceso \cite{arinde2024multisensor}.
  \vspace{3mm}
  \begin{columns}[T]
    \column{0.32\textwidth}
    \begin{block}{Desconexión}
      Nadie relaciona una apertura con un golpe de hace 5 segundos. Se pierden los primeros segundos.
    \end{block}
    \column{0.32\textwidth}
    \begin{block}{Costo}
      Lo comercial es caro. Lo barato trae un solo sensor y rara vez avisa a distancia \cite{oktafiani2025esp32}.
    \end{block}
    \column{0.32\textwidth}
    \begin{block}{Sin métricas}
      El dueño no sabe cuántas alertas fueron falsas ni cuánto tardó un aviso \cite{pugh2015fusion}.
    \end{block}
  \end{columns}
  \vspace{3mm}
  \begin{alertblock}{Problema}
    Falta un sistema integrado y barato que correlacione varios sensores, responda sin depender de Internet y mida su propio desempeño.
  \end{alertblock}
\end{frame}

% ===============================================================
\section{Marco teórico}

\begin{frame}{Marco teórico: arquitectura y sensores}
  \begin{columns}[T]
    \column{0.47\textwidth}
    \textbf{Arquitectura IoT por capas} \cite{alfuqaha2015survey, ray2018survey}\par\vspace{3mm}
    \begin{tikzpicture}[node distance=2mm]
      \node[capa=cApp] (ap) {Aplicación};
      \node[capa=cProc, below=of ap] (rd) {Red};
      \node[capa=cSensor, below=of rd] (pe) {Percepción};
      \node[nuestro=cApp, right=3mm of ap] {App web en EC2};
      \node[nuestro=cProc, right=3mm of rd] {Wi-Fi + MQTT/TLS};
      \node[nuestro=cSensor, right=3mm of pe] {Sensores y actuadores};
    \end{tikzpicture}
    \column{0.5\textwidth}
    \textbf{Sensores}\par\vspace{2mm}
    \begin{columns}[T]
      \column{0.32\linewidth}\foto{pir.jpg}{2.4cm}{PIR HC-SR501}{Calor infrarrojo en movimiento, 3 a 7\,m}
      \column{0.32\linewidth}\foto{reed.jpg}{2.4cm}{Magnético MC-38}{Imán y reed switch, detecta apertura}
      \column{0.32\linewidth}\foto{sw420.jpg}{2.4cm}{Vibración SW-420}{Golpes, comparador LM393}
    \end{columns}\par\vspace{2mm}
    \footnotesize El PIR aparece en 7 de 10 sistemas ESP32 revisados \cite{oktafiani2025esp32}.
  \end{columns}
  \credito{Fotos referenciales de Wikimedia Commons: PIR por Lethalattraction (CC BY-SA 4.0), reed switch por Satishbhandare (dominio público), SW-420 por Suyash Dwivedi (CC BY-SA 4.0).}
\end{frame}

\begin{frame}{Marco teórico: MQTT, borde y nube}
  \begin{columns}[c]
    \column{0.56\textwidth}
    \centering
    \begin{tikzpicture}[font=\small]
      \node[n=cProc] (esp) at (0,0) {ESP32\\{\scriptsize publica}};
      \node[n=cApp]  (brk) at (3.6,0) {Mosquitto\\{\scriptsize broker}};
      \node[n=cApp]  (bk)  at (6.3,0.85) {Backend\\{\scriptsize suscrito}};
      \node[n=cApp]  (web) at (6.3,-0.85){App web};
      \draw[flecha] (esp) -- node[above=1pt, font=\tiny\ttfamily] {casa/eventos} (brk);
      \draw[flecha] (brk.east) -- (bk.west);
      \draw[flecha] (bk) -- (web);
      \node[font=\scriptsize, text=gris, below=2mm of esp] {decide y actúa};
      \node[font=\scriptsize, text=gris, below=2mm of brk] {TLS, puerto 8883};
    \end{tikzpicture}
    \column{0.42\textwidth}
    \begin{itemize}
      \item MQTT usa publicación y suscripción, con una cabecera mínima de 2 bytes \cite{alfuqaha2015survey, light2017mosquitto}.
      \item Las decisiones urgentes se toman en el borde, sin esperar a Internet \cite{shi2016edge}.
      \item La nube guarda los eventos y calcula métricas.
      \item Correlacionar varios sensores ayuda a controlar las falsas alarmas \cite{walchli2009building, pugh2015fusion}.
    \end{itemize}
  \end{columns}
\end{frame}

% ===============================================================
\section{Estado del arte}

\begin{frame}{Estado del arte}
  \footnotesize
  \renewcommand{\arraystretch}{1.15}
  \begin{tabular}{@{}l>{\raggedright\arraybackslash}p{7.2cm}>{\raggedright\arraybackslash}p{4.9cm}@{}}
    \toprule
    \textbf{Integrante} & \textbf{Artículo} & \textbf{Qué tomamos} \\
    \midrule
    Ariann & Nannung et al. \cite{nannung2025iot}. ESP32-CAM con reconocimiento facial (95,31\%), DHT22 y MQ-2 & Relé con cerradura solenoide \\
     & Oktafiani y Widhiantoro \cite{oktafiani2025esp32}. Revisión de 10 sistemas ESP32 & PIR como sensor principal \\
    \midrule
    Diego & Wälchli y Braun \cite{walchli2009building}. Oficina con 10 nodos, hasta 90\% menos comunicación & Reportar eventos, no lecturas \\
     & Pugh et al. \cite{pugh2015fusion}. Fusión de varios PIR con límite de falsas alarmas & Regla de dos sensores en 10\,s \\
    \midrule
    Federico & Ray \cite{ray2018survey}. Arquitecturas IoT ordenadas por dominio & Modelo por capas \\
     & Al-Fuqaha et al. \cite{alfuqaha2015survey}. Tecnologías y protocolos IoT & Elección de MQTT \\
    \midrule
    Sofia & Jin et al. \cite{jin2018wlid}. Intrusos detectados por RSSI de Wi-Fi, cerca de 100\% & RSSI como métrica del enlace \\
     & Arinde e Idowu \cite{arinde2024multisensor}. PIR + magnético + aviso por GSM & La misma base de sensores, con Wi-Fi \\
    \bottomrule
  \end{tabular}
\end{frame}

% ===============================================================
\section{Metodología}

\begin{frame}{Componentes y microcontrolador}
  \begin{columns}[T]
    \column{0.5\textwidth}
    \begin{block}{Núcleo (cumple los mínimos)}
      \small
      \begin{itemize}
        \item 3 sensores: PIR, magnético MC-38 y SW-420
        \item Actuadores: sirena 12\,V y cerradura solenoide
        \item Inalámbrico: Wi-Fi 2,4\,GHz + MQTT/TLS
        \item Interfaz: app web propia con métricas
        \item Enclosure: carcasa impresa en 3D (PLA)
        \item Relés de 2 canales, fuente de 12\,V y 2\,A
      \end{itemize}
    \end{block}
    \column{0.46\textwidth}
    \begin{block}{ESP32 DevKit V1}
      \small
      \begin{itemize}
        \item Doble núcleo a 240\,MHz, 520\,KB de SRAM
        \item Wi-Fi 802.11 b/g/n y Bluetooth 4.2
        \item Lógica de 3,3\,V, entradas de sobra
        \item Hora por NTP para medir la latencia
        \item Barato, como se recomendó
      \end{itemize}
    \end{block}
  \end{columns}
\end{frame}

\begin{frame}{Hardware de referencia}
  \begin{columns}[T]
    \column{0.19\textwidth}\foto{esp32.jpg}{2.5cm}{ESP32 DevKit V1}{Procesa y publica}
    \column{0.19\textwidth}\foto{rele.jpg}{2.5cm}{Módulo de relés}{Conmuta cargas de 12\,V}
    \column{0.19\textwidth}\foto{sirena.jpg}{2.5cm}{Sirena 12\,V}{Alerta sonora}
    \column{0.19\textwidth}\foto{cerradura.jpg}{2.5cm}{Cerradura 12\,V}{Bloqueo de la puerta}
    \column{0.19\textwidth}\foto{esp32cam.jpg}{2.5cm}{ESP32-CAM}{Opcional}
  \end{columns}
  \vspace{4mm}
  \begin{block}{Opcionales, si alcanza el tiempo}
    \small ESP32-CAM (OV2640, 2\,MP) que fotografía cada evento crítico y la sube al servidor. Lector de huella R307 para desarmar en el sitio. Sensor de humo MQ-2.
  \end{block}
  \credito{Fotos referenciales de Wikimedia Commons: ESP32 por Edwiyanto, relé por Suyash Dwivedi, ESP32-CAM por Nowforever (CC BY-SA 4.0), sirena por Harrihealey02 (CC BY-SA 3.0), cerradura electromagnética por Mooinglemur (dominio público).}
\end{frame}

\begin{frame}{Diagrama de bloques}
  \begin{columns}[c]
    \column{0.62\textwidth}
    \centering{\let\footnotesize\scriptsize\let\scriptsize\tiny
      \resizebox{!}{0.8\textheight}{\begin{tikzpicture}[font=\footnotesize]
  % Capa de percepción
  \node[caja=cSensor] (pir) at (-2.65,0)  {Sensor PIR\\{\scriptsize HC-SR501}};
  \node[caja=cSensor] (mag) at (0,0)      {Sensor magnético\\{\scriptsize MC-38}};
  \node[caja=cSensor] (vib) at (2.65,0)   {Sensor de vibración\\{\scriptsize SW-420}};
  % Procesamiento
  \node[caja=cProc, minimum width=2.2cm, minimum height=1cm] (esp) at (0,-1.9)
    {\textbf{Microcontrolador}\\\textbf{ESP32}};
  % Nube e interfaz
  \node[caja=cApp, minimum width=1.8cm] (srv) at (3.0,-1.9) {Servidor AWS\\{\scriptsize MQTT + backend}};
  \node[caja=cApp, minimum width=1.8cm] (web) at (3.0,-3.6) {App web\\{\scriptsize métricas}};
  % Actuación
  \node[caja=cPower, minimum width=2cm] (psu) at (-2.75,-3.6) {Fuente\\externa 12\,V};
  \node[caja=cAct] (relay) at (0,-3.6)    {Módulo de relés};
  \node[caja=cAct] (sirena) at (-1.3,-5.2) {Sirena\\{\scriptsize 12\,V}};
  \node[caja=cAct] (lock) at (1.3,-5.2)    {Cerradura\\solenoide {\scriptsize 12\,V}};

  % Sensores -> ESP32 (entradas GPIO)
  \draw[flecha] (pir.south) -- ++(0,-0.45) -| ([xshift=-7mm]esp.north);
  \draw[flecha] (mag.south) -- (esp.north)
    node[etiqueta, pos=0.4] {GPIO};
  \draw[flecha] (vib.south) -- ++(0,-0.45) -| ([xshift=7mm]esp.north);

  % ESP32 -> nube -> app
  \draw[flecha, dashed, draw=cApp] (esp.east) -- (srv.west)
    node[midway, above, font=\scriptsize, text=cApp] {Wi-Fi};
  \node[font=\scriptsize, text=cApp, below=0pt of $(esp.east)!0.5!(srv.west)$] {MQTT};
  \draw[flecha, draw=cApp] (srv.south) -- (web.north)
    node[midway, left, font=\scriptsize, text=cApp] {HTTPS};

  % ESP32 -> relés -> actuadores
  \draw[flecha] (esp.south) -- (relay.north)
    node[etiqueta, midway] {GPIO};
  \draw[flecha, draw=cPower] (psu.east) -- (relay.west)
    node[midway, above, font=\scriptsize, text=cPower] {12\,V};
  \draw[flecha] (relay.south) -- ++(0,-0.4) -| (sirena.north);
  \draw[flecha] (relay.south) -- ++(0,-0.4) -| (lock.north);
\end{tikzpicture}
}}
    \column{0.34\textwidth}
    \small
    \begin{itemize}
      \item {\color{cSensor}\textbf{Azul}}: adquisición
      \item {\color{cProc}\textbf{Naranja}}: procesamiento local
      \item {\color{cAct}\textbf{Verde}}: actuación
      \item {\color{cApp}\textbf{Morado}}: nube e interfaz
    \end{itemize}
    \vspace{3mm}
    {\footnotesize Los actuadores usan la fuente de 12\,V, nunca el ESP32.}
  \end{columns}
\end{frame}

\begin{frame}{Lógica de detección}
  \begin{center}
  \begin{tikzpicture}[font=\footnotesize]
    \node[f=cSensor] (leer) at (0,0)     {Leer\\sensores};
    \node[d]         (d1)   at (2.35,0)  {¿Sensor\\activo?};
    \node[f=cSensor] (cor)  at (4.75,0)  {Correlacionar\\últimos 10\,s};
    \node[d]         (d2)   at (7.15,0)  {¿Dos\\sensores?};
    \node[f=cAct]    (act)  at (9.55,0)  {Sirena y\\bloqueo};
    \node[f=cApp]    (pub)  at (11.9,0)  {Publicar\\por MQTT};
    \draw[flecha] (leer) -- (d1);
    \draw[flecha] (d1) -- node[above, font=\scriptsize] {Sí} (cor);
    \draw[flecha] (cor) -- (d2);
    \draw[flecha] (d2) -- node[above, font=\scriptsize] {Sí} (act);
    \draw[flecha] (act) -- (pub);
    \draw[flecha] (d1.north) -- ++(0,0.45) -| node[pos=0.25, above, font=\scriptsize] {No} (leer.north);
    \draw[flecha, draw=cProc] (d2.south) -- ++(0,-0.4) -| node[pos=0.25, below, font=\scriptsize, text=cProc] {No: alerta leve} (pub.south);
    \draw[flecha, draw=gris] (pub.north) -- ++(0,0.65) -| node[pos=0.25, above, font=\scriptsize, text=gris] {vuelve a leer} ([xshift=-3mm]leer.north);
  \end{tikzpicture}
  \end{center}
  \begin{itemize}
    \item \textbf{Alerta leve}: un solo sensor. Se publica el evento, sin sirena.
    \item \textbf{Evento crítico}: dos sensores distintos en menos de 10\,s, por ejemplo vibración y luego apertura. Suenan la sirena y el bloqueo.
    \item Cada mensaje lleva sensor, nivel, número de secuencia y hora. Cada 30\,s sale un latido con el RSSI.
  \end{itemize}
\end{frame}

\begin{frame}{Interfaz con procesamiento previo}
  \begin{columns}[T]
    \column{0.36\textwidth}
    \begin{block}{AWS EC2 (t3.micro)}
      \small
      \begin{itemize}
        \item Mosquitto (MQTT/TLS)
        \item Backend en FastAPI
        \item Base SQLite
        \item App web
      \end{itemize}
      \footnotesize Todo en Docker Compose, así corre igual en una laptop.
    \end{block}
    \column{0.6\textwidth}
    \small
    \renewcommand{\arraystretch}{1.25}
    \begin{tabular}{@{}l>{\raggedright\arraybackslash}p{5.2cm}@{}}
      \toprule
      \textbf{Métrica} & \textbf{Cálculo} \\
      \midrule
      Latencia & Recepción menos hora del evento. Promedio y p95 \\
      Pérdida & Huecos en el número de secuencia \\
      Disponibilidad & Latidos recibidos sobre los esperados \\
      Eventos & Conteo por sensor y por hora \\
      Falsas alarmas & Alertas que el usuario marca como falsas \\
      Enlace & RSSI promedio en dBm \\
      \bottomrule
    \end{tabular}
  \end{columns}
\end{frame}

% ===============================================================
\section{Objetivos y alcances}

\begin{frame}{Objetivos}
  \begin{block}{Objetivo general}
    Diseñar e implementar un sistema inteligente de seguridad residencial IoT que detecte intrusiones con tres sensores, clasifique el riesgo en un ESP32, responda con sirena y bloqueo, y avise por Wi-Fi a una app web que mida su propio desempeño.
  \end{block}
  \textbf{Objetivos específicos}
  \begin{columns}[T]
    \column{0.48\textwidth}
    \begin{itemize}
      \item Detectar movimiento, apertura y vibración
      \item Clasificar cada evento como leve o crítico con la regla de 10\,s
      \item Activar sirena y cerradura sin depender de Internet
    \end{itemize}
    \column{0.48\textwidth}
    \begin{itemize}
      \item Alertas en menos de 2\,s (p95) vía MQTT a AWS
      \item App web con 6 métricas de desempeño
      \item Hardware dentro de una carcasa 3D
      \item (Opcional) Huella y ESP32-CAM
    \end{itemize}
  \end{columns}
\end{frame}

\begin{frame}{Alcances y limitaciones}
  \begin{columns}[T]
    \column{0.45\textwidth}
    \begin{block}{Alcances}
      \small
      Una zona de prueba con una puerta, un PIR y un sensor de vibración. Incluye un nodo ESP32, sirena, cerradura, servidor en AWS y app web.\par\vspace{2mm}
      Huella, cámara y humo quedan como opcionales.
    \end{block}
    \column{0.51\textwidth}
    \begin{alertblock}{Limitaciones}
      \small
      \begin{itemize}
        \item Un solo nodo, PIR hasta 7\,m y Wi-Fi de 2,4\,GHz
        \item Sin Internet la sirena sigue, pero la app no recibe eventos
        \item Sin batería de respaldo
        \item No cubre sabotaje del cableado ni interferencia de radio
        \item Si el laboratorio bloquea el puerto 8883, usamos un hotspot
      \end{itemize}
    \end{alertblock}
  \end{columns}
\end{frame}

% ===============================================================
\begin{frame}{Referencias}
  \scriptsize
  \begin{thebibliography}{10}
  \bibitem{alfuqaha2015survey} A. Al-Fuqaha et al., ``Internet of Things: A survey on enabling technologies, protocols, and applications,'' \emph{IEEE Commun. Surveys Tuts.}, vol. 17, no. 4, pp. 2347--2376, 2015.
  \bibitem{oktafiani2025esp32} F. Oktafiani y D. Widhiantoro, ``Studi literatur penggunaan ESP32 untuk sistem keamanan lingkungan rumah,'' \emph{SNIV}, vol. 4, no. 1, pp. 267--274, 2025.
  \bibitem{arinde2024multisensor} V. Arinde y L. Idowu, ``Multi-sensor intrusion detection system,'' arXiv:2406.05137, 2024.
  \bibitem{pugh2015fusion} M. Pugh, J. Brewer y J. Kvam, ``Sensor fusion for intrusion detection under false alarm constraints,'' \emph{IEEE SAS}, pp. 1--6, 2015.
  \bibitem{ray2018survey} P. P. Ray, ``A survey on Internet of Things architectures,'' \emph{J. King Saud Univ. Comput. Inf. Sci.}, vol. 30, no. 3, pp. 291--319, 2018.
  \bibitem{nannung2025iot} J. Nannung et al., ``Development of an IoT-based home security system prototype using multisensors and ESP32-CAM,'' \emph{JESSI}, vol. 6, no. 2, pp. 168--177, 2025.
  \bibitem{light2017mosquitto} R. A. Light, ``Mosquitto: server and client implementation of the MQTT protocol,'' \emph{JOSS}, vol. 2, no. 13, p. 265, 2017.
  \bibitem{shi2016edge} W. Shi et al., ``Edge computing: Vision and challenges,'' \emph{IEEE Internet Things J.}, vol. 3, no. 5, pp. 637--646, 2016.
  \bibitem{walchli2009building} M. Wälchli y T. Braun, ``Building intrusion detection with a wireless sensor network,'' \emph{ADHOCNETS 2009}, LNICST 28, pp. 607--622, 2010.
  \bibitem{jin2018wlid} Y. Jin et al., ``A whole-home level intrusion detection system using WiFi-enabled IoT,'' \emph{IWCMC}, pp. 494--499, 2018.
  \end{thebibliography}
\end{frame}

\end{document}
